\documentclass[11pt,a4paper]{article}
\usepackage[T1]{fontenc}
\usepackage[utf8]{inputenc}
\usepackage[french]{babel}
\usepackage{amsmath,amssymb}
\usepackage[margin=2.5cm]{geometry}
\usepackage{pgfplots}
\pgfplotsset{compat=1.18}
\usepackage{hyperref}
\hypersetup{hidelinks}
\definecolor{partiA}{HTML}{FE0000}
\definecolor{partiB}{HTML}{F15F24}
\definecolor{partiC}{HTML}{FEB616}
\definecolor{partiD}{HTML}{FFE500}
\pgfplotsset{pouls/.style={width=0.96\linewidth,height=5.1cm,xmin=0,xmax=500,
  xlabel={Tour},grid=major,legend style={font=\small},
  cycle list={{partiA,thick},{partiB,thick},{partiC,thick},{partiD,thick}}}}

\title{Pouls\\L'espace de la politique\\
\large Compétition spatiale, abstention, réputation et renouvellement}
\author{}
\date{}
\begin{document}
\maketitle
\begin{abstract}
Pouls est un modèle exploratoire du marché politique comme marché du
<< moins pire >>, organisé sur un axe gauche--droite. Il associe un choix probabiliste entre offres acceptables,
un seuil d'abstention, des coûts de réputation endogènes et l'entrée de
nouveaux candidats dans les zones délaissées. Cette note décrit le moteur
effectivement utilisé et rapporte neuf simulations reproductibles de
500 tours. Il ne s'agit ni d'une estimation sur données électorales ni
d'une prévision.
\end{abstract}

\section{Références théoriques}

\subsection{Le marché politique comme marché du moins pire}
Le raisonnement repose sur l'énoncé suivant :
\begin{quote}
Le marché politique est un marché linéaire allant de la droite à la
gauche. C'est un marché du moins pire : on vote pour ce qui nous
convient le moins mal. Les effets de réputation sont déterminants
et pèsent sur les coûts du vote.
\end{quote}
Cet énoncé constitue le parti pris théorique du modèle, et non une loi
empirique établie par les simulations. Le terme \emph{linéaire} décrit
ici l'organisation des offres sur un axe politique unique : les
positions se comparent et s'ordonnent de la droite à la gauche.
Il ne signifie pas que le coût de la distance doit être une fonction
linéaire. Les coûts linéaire et quadratique sont tous deux compatibles
avec ce même espace politique unidimensionnel. L'orientation graphique
retenue, de $-1$ à gauche à $+1$ à droite, ne change pas ce raisonnement.

Parler de marché revient à mettre en relation une distribution de
préférences électorales et des offres politiques concurrentes. Chaque
électeur a une position qui lui conviendrait idéalement ; les partis
proposent des positions qui ne coïncident généralement pas exactement
avec cette préférence. Il ne s'agit pas d'un marché à prix monétaire :
le coût représente le désaccord et la difficulté à accepter une offre,
non une somme effectivement payée.

Le vote est alors envisagé comme un choix sous contrainte entre des
offres imparfaites, plutôt que comme l'adhésion nécessaire à une offre
idéale. On ne vote pas forcément pour ce qui convient pleinement,
mais pour ce qui convient \emph{le moins mal} parmi les possibilités
disponibles. Ce raisonnement distingue deux questions : quelle offre
est relativement moins coûteuse que les autres, et existe-t-il encore
une offre dont le coût reste acceptable ? La première organise la
concurrence entre partis ; la seconde ouvre la possibilité de
l'abstention. Même la meilleure offre disponible peut être trop
coûteuse pour justifier un vote.

\subsection{Différenciation spatiale}
Hotelling \cite{hotelling1929} fournit le cadre de la concurrence entre
offres localisées sur un axe, avec un coût lié à leur éloignement des
consommateurs. Ici, les offres sont des candidats et la distance représente
un désaccord politique. Le repère quadratique vient de
d'Aspremont, Gabszewicz et Thisse \cite{daspremont1979}.
La référence est bien de 1979. Ces articles concernent des modèles de
concurrence économique ; leurs résultats d'équilibre ne se transposent
pas automatiquement au présent jeu politique sans prix.

\subsection{Aliénation et limite du choix du moins pire}
Le coût maximal accepté a une justification politique précise :
\emph{l'électeur peut refuser de voter même pour le candidat le plus
proche, parce que celui-ci est encore trop éloigné de ses préférences}.
Il ne suffit donc pas d'être le moins éloigné des candidats pour
constituer une offre acceptable. C'est l'abstention par
\emph{aliénation}, étudiée dans les modèles spatiaux de participation,
notamment par Hinich et Ordeshook \cite{hinich1969}.
Adams, Dow et Merrill \cite{adams2006} distinguent explicitement
ses conséquences de celles de l'abstention par indifférence.

L'indifférence concerne le faible écart entre les offres : l'électeur
les juge trop semblables pour que leur départage justifie sa
participation. L'aliénation concerne au contraire leur niveau absolu
d'acceptabilité : aucune offre ne lui convient suffisamment, même
s'il préfère nettement l'une aux autres. Un électeur peut donc être
aliéné sans être indifférent. Dans Pouls, seule cette seconde logique
déclenche l'abstention ; la proximité des coûts entre candidats ne
provoque pas, à elle seule, de retrait du vote.

Le seuil $k_{\max}$ représente ainsi une limite au compromis politique
que l'électeur accepte de consentir, et non un simple coût matériel
pour se rendre aux urnes. La règle est
\[
\text{abstention de }x\text{ au tour }t
\quad\Longleftrightarrow\quad
\min_i c_{i,t}(x)>k_{\max}.
\]
Le << moins pire >> reste un choix relatif, mais il n'oblige pas à
accepter n'importe quelle offre : lorsque même la moins coûteuse
dépasse cette limite, l'électeur s'abstient.

Dans une version purement spatiale, à coûts de base identiques et sans
bruit, la moins coûteuse est aussi la plus proche. Pouls étend cette
logique au coût total, qui inclut la réputation et les chocs
individuels : le candidat géométriquement le plus proche n'est donc
pas nécessairement le moins coûteux. Le seuil commun à tous les
électeurs et sa forme de coupure nette sont des hypothèses
simplificatrices du modèle, non des résultats établis par ces
références.

\subsection{La réputation comme composante déterminante du coût}
La proximité politique ne suffit pas à caractériser ce qui convient
le moins mal. Deux partis situés au même endroit peuvent présenter
des coûts différents si leur réputation diffère. Dans Pouls, cette
réputation agit sur le coût de base $k_{i,t}$, qui s'ajoute au coût
de distance : elle pèse donc sur l'ensemble de l'offre, y compris
pour un électeur situé exactement à la position du parti.
Une offre plus éloignée peut ainsi être préférée à une offre plus
proche si son coût de base est suffisamment inférieur.

La réputation ne désigne pas ici une qualité objective mesurée
indépendamment du vote. Elle est représentée, par hypothèse, par le
succès électoral récent. Un parti qui dépasse la part égale $1/N$
voit son coût diminuer ; un parti qui reste en dessous voit son coût
augmenter. Le soutien passé modifie ainsi les conditions du choix
présent. La concurrence est cumulative : les partis n'affrontent pas
seulement une distribution de préférences, mais aussi l'histoire de
leurs succès et de leurs échecs.

Les effets de réputation sont déterminants en un double sens dans
ce mécanisme : ils modifient les probabilités relatives de vote
\emph{et} peuvent faire entrer une offre dans l'ensemble acceptable
ou l'en faire sortir. Le paramètre $\eta$ règle leur intensité ;
leur poids effectif dépend des paramètres et de la trajectoire.
Il ne s'agit donc pas d'affirmer qu'ils dominent toujours la distance.
Le seuil d'abstention et le renouvellement des partis permettent
d'étudier comment cet avantage acquis rencontre ses limites.

\subsection{Choix probabiliste et extensions}
Le choix discret probabiliste s'inscrit dans la tradition présentée par
Manski et McFadden \cite{manski1981}. Le softmax employé est un logit
multinomial de température fixée à 1, conditionnel à l'ensemble des offres
acceptables. L'abstention n'est pas une option extérieure du logit :
elle intervient seulement si cet ensemble est vide.
Le bruit gaussien additif du coût est une hypothèse supplémentaire,
distincte des erreurs de type valeur extrême qui peuvent fonder un logit.
Le softmax ne découle donc pas ici du seul bruit gaussien.

Le << moins pire >> est ainsi traduit de façon probabiliste, et non
par une règle qui attribuerait toutes les voix à l'offre la moins
coûteuse. Pour deux offres admissibles $i$ et $j$, leurs chances
relatives vérifient
\[
\frac{P_{i,t}(x)}{P_{j,t}(x)}
=\exp\!\left[c_{j,t}(x)-c_{i,t}(x)\right].
\]
Un coût plus faible rend donc une offre plus probable sans rendre
son choix certain. Si toutes les offres dépassent le coût maximal
accepté, aucune comparaison relative ne suffit à provoquer le vote :
l'électeur s'abstient.

La réputation à mémoire courte, la mobilité bornée et le renouvellement
des candidats sont des hypothèses propres à Pouls, et non des résultats
attribués aux références précédentes.

\section{Hypothèses du modèle}

\subsection{Électorat et conditions initiales}
Les $N$ places de candidature, $2\leq N\leq7$, ont des positions
$p_{i,t}\in[-1,1]$. L'électorat fixe a une densité $f$ de masse 1.
L'application propose $f(x)=1/2$, ou le mélange
\[
f(x)=w h(x;\mu_L,\sigma_L)+(1-w)h(x;\mu_R,\sigma_R),
\quad
h(x;\mu,\sigma)=
\frac{\phi((x-\mu)/\sigma)}
{\sigma[\Phi((1-\mu)/\sigma)-\Phi((-1-\mu)/\sigma)]}.
\]
Chaque normale est tronquée et normalisée séparément sur $[-1,1]$.
Le poids $w$ décrit une composante, pas la proportion située à gauche
de zéro ; le mélange n'a pas nécessairement deux modes.
Au chargement, l'électorat est uniforme et quatre candidats commencent
en $p_{i,0}=0$, avec $k_{i,0}=k_0=1$.
Le vote initial fournit la première observation de réputation.

\subsection{Coût additif, softmax et abstention}
Pour l'électeur $x$, le coût de l'offre $i$ au tour $t$ est
\begin{equation}\label{eq:cost}
\boxed{c_{i,t}(x)=\max\{0,\ k_{i,t}+\beta|x-p_{i,t}|^\alpha+
\varepsilon_{i,t}(x)\}},\quad
\varepsilon_{i,t}(x)\sim\mathcal N(0,\sigma^2).
\end{equation}
On impose $\alpha>0$, $\beta\geq0$, $k_{i,t}>0$ et $\sigma\geq0$.
Le bruit est \emph{additif}, jamais multiplicatif ; $\sigma=0{,}1$
désigne une unité de coût, pas 10\,\%. Un choc par point d'électorat et
place de candidat est renouvelé à chaque tour, puis fixé pour toutes
les positions testées, les décisions et le vote de ce tour.

Le seuil $k_{\max}$ est le coût total maximal accepté, et non une borne
portant seulement sur $k_i$. L'ensemble admissible dépend de l'électeur :
\[
\mathcal E_t(x)=\{i:c_{i,t}(x)\leq k_{\max}\}.
\]
La probabilité de vote est
\begin{equation}
P_{i,t}(x)=
\begin{cases}
\displaystyle\frac{\exp[-c_{i,t}(x)]}
{\sum_{j\in\mathcal E_t(x)}\exp[-c_{j,t}(x)]},
& i\in\mathcal E_t(x),\\[6pt]
0,&\text{sinon}.
\end{cases}
\end{equation}
Si $\mathcal E_t(x)=\varnothing$, l'électeur s'abstient.
L'égalité au seuil est acceptée. À coûts finis, tout candidat admissible
a une probabilité strictement positive : le choix n'est pas déterministe.
L'intégration utilise les probabilités, sans simuler un bulletin individuel.

Pour expliciter le lien avec l'aliénation spatiale, supposons
$k_{i,t}=k$ pour tous les candidats, $\sigma=0$, $\beta>0$ et
$k\leq k_{\max}$. Une offre est acceptable si et seulement si
\[
|x-p_{i,t}|\leq
r,\qquad
r=\left(\frac{k_{\max}-k}{\beta}\right)^{1/\alpha}.
\]
L'électeur s'abstient donc lorsque le candidat le plus proche est
encore à une distance strictement supérieure à $r$. Si $k>k_{\max}$,
aucune offre n'est acceptable, même à distance nulle. Avec des
réputations différentes ou du bruit, c'est la condition sur le coût
total, et non ce rayon commun, qui s'applique.

Les voix sur l'ensemble de l'électorat, la participation et les parts
des exprimés sont
\[
v_{i,t}=\int P_{i,t}(x)f(x)\,dx,\qquad
T_t=\sum_i v_{i,t},\qquad
s_{i,t}=\begin{cases}v_{i,t}/T_t,&T_t>0,\\0,&T_t=0.\end{cases}
\]
L'abstention vaut $a_t=1-T_t$. Numériquement, les exprimés sont mis à
zéro si $T_t\leq10^{-10}$, mais un renouvellement exige toujours des
\emph{voix réellement nulles}, pas cette convention ni un arrondi d'affichage.
Les parts principales de l'application sont les $v_i$ ; la réputation
emploie les $s_i$. Ainsi des parts totalisant 100\,\% des électeurs
ne peuvent pas coexister avec 30\,\% d'abstention.

\subsection{Réputation et mémoire individuelle}
Soit $h_{i,t-1}\in\{0,1,2,3\}$ le nombre d'observations disponibles
pour l'occupant actuel de la place $i$. Son succès passé est
\[
S_{i,t-1}=
\begin{cases}
\displaystyle\frac{\sum_{\ell=0}^{h_{i,t-1}-1}(3-\ell)s_{i,t-1-\ell}}
{\sum_{\ell=0}^{h_{i,t-1}-1}(3-\ell)},&h_{i,t-1}>0,\\[6pt]
1/N,&h_{i,t-1}=0.
\end{cases}
\]
Pour un candidat qui reste en lice, le coût est actualisé avant le vote :
\begin{equation}\label{eq:reputation}
\boxed{k_{i,t}=k_{i,t-1}\exp[\eta(1/N-S_{i,t-1})]},\qquad\eta\geq0.
\end{equation}
Les trois poids sont donc $3,2,1$, avec renormalisation au démarrage.
Il n'y a ni lissage exponentiel des scores ni mémoire des déplacements.
Un succès supérieur à $1/N$ réduit le coût ; un succès inférieur
l'augmente. À $\eta=0$, les coûts restent fixes \emph{entre renouvellements}.
Le score produit au tour $t$ n'intervient qu'au tour suivant.
Les coûts ne changent pas pendant les recherches de position.

\subsection{Innovation : remplacer une offre sans voix}
Au début du tour $t$, une place telle que $v_{i,t-1}=0$ accueille un
nouvel entrant, sous réserve qu'il existe des abstentionnistes.
Le coût est remis à \textbf{1}, même si $k_0\neq1$, et la mémoire du
prédécesseur est oubliée pour cette seule place. Le nouveau candidat
ne subit pas l'actualisation de réputation de son prédécesseur.

Soit $\mathcal A_{t-1}=\{x:\mathcal E_{t-1}(x)=\varnothing\}$ le groupe
des abstentionnistes, reconstitué avec les positions, coûts et chocs du
\emph{vote précédent}. Pour une arrivée en $q$, on maximise
\begin{equation}\label{eq:entry}
R_{i,t}(q)=\int_{\mathcal A_{t-1}}
\mathbf1\!\left\{\max(0,1+\beta|x-q|^\alpha+
\varepsilon_{i,t}(x))\leq k_{\max}\right\}f(x)\,dx.
\end{equation}
La position d'entrée maximise donc la \textbf{masse d'abstentionnistes
accessibles}, et non une densité ponctuelle, un tirage pondéré ou les
voix totales anticipées. La recherche parcourt tout l'axe au pas $0{,}01$,
avec les chocs du nouveau tour. Le hasard départage seulement les maxima
égaux. Sans bruit et pour $\beta>0$, $k_{\max}\geq1$, la zone couverte
est de rayon $[(k_{\max}-1)/\beta]^{1/\alpha}$ autour de $q$.

Les arrivées multiples ciblent indépendamment le même groupe passé :
elles peuvent donc se concurrencer. La couverture n'est pas un nombre
de votes garanti, puisque le vote reste probabiliste et que les autres
offres changent. Si aucune offre nouvelle ne peut rendre ces électeurs
admissibles, le maximum peut être nul ; l'entrée n'assure pas le succès.
Sans abstentionnistes, le renouvellement est reporté et signalé.
Les lettres et familles de couleurs identifient les places. Dans
l'application et la vidéo, chaque génération de parti reçoit une
nouvelle teinte ; les courbes sont interrompues à la naissance du
nouveau parti, plutôt que de suggérer un déplacement du prédécesseur.

\subsection{Mobilité et calendrier d'un tour}
La première position d'un entrant n'est pas soumise à $\gamma$.
Il ne se déplace pas une seconde fois lors de son arrivée. Les autres
candidats observent les arrivées, puis choisissent simultanément, chacun
contre les positions de ce même état intermédiaire $\widetilde{\mathbf p}_t$ :
\[
p_{i,t}\in\operatorname*{arg\,max}_{q\in Q_{i,t}}
\mathcal V_i(q,\widetilde{\mathbf p}_{-i,t};
             \mathbf k_t,\boldsymbol\varepsilon_t).
\]
Ici $\mathcal V_i$ est la masse de voix espérées sur \emph{tout}
l'électorat. L'ensemble $Q_{i,t}$ contient la grille de pas $0{,}01$
dans $[p_{i,t-1}-\gamma,p_{i,t-1}+\gamma]\cap[-1,1]$,
la position courante et les deux bornes exactes de cet intervalle.
Ainsi $|p_{i,t}-p_{i,t-1}|\leq\gamma$ pour les candidats en place.
$\gamma=0$ les immobilise ; $\gamma=2$ ouvre tout l'axe.
Sans gain supérieur à $10^{-10}$, la position courante est conservée.
Pour ces meilleures réponses ordinaires, le parcours gauche--droite
peut départager des solutions nouvelles équivalentes.

Un tour suit donc : (1) reconstruction du groupe abstentionniste et
arrivées ; (2) actualisation de réputation des occupants conservés ;
(3) meilleures réponses simultanées de ces derniers ; (4) application
collective, vote et ajout des scores aux mémoires individuelles.
Une optimisation individuelle ne renouvelle ou ne déplace que la place
sélectionnée, mais actualise tous les coûts et toutes les mémoires.
Les gains anticipés ne sont pas nécessairement réalisés quand les
concurrents changent aussi de position.

\subsection{Paramètres et implémentation}
Les valeurs par défaut sont
\[
N=4,\quad\alpha=1,\quad\beta=3{,}5,\quad k_0=1,\quad
k_{\max}=1{,}25,\quad\gamma=0{,}1,\quad\eta=0{,}04,\quad\sigma=0{,}1.
\]
L'intégration emploie $1\,600$ points médians et leur masse normalisée.
Le softmax est stabilisé en retranchant le coût admissible minimal.
Chaque simulation contient au plus 500 tours, soit 501 états avec le
départ. Les graphes sont actualisés pas à pas ; les points cerclés
signalent les arrivées. L'abstention est blanche et les sept couleurs
de candidature sont contrastées. Les diagrammes de phase sont supprimés.
Les graines, coûts, scores, arrivées et longueurs de mémoire sont
exportés ; une interruption conserve les tours achevés.

\section{Résultats des simulations}

\subsection{Protocole reproductible}
Les résultats ci-dessous ont été calculés avec le moteur commun
\texttt{model.js}, sans ajustement forcé des parts de voix.
Trois scénarios sont comparés, chacun pour les graines maîtresses
$314159$, $271828$ et $161803$ : la référence aux paramètres ci-dessus,
la suppression de la réputation ($\eta=0$), et le coût quadratique
($\alpha=2$). Les autres paramètres restent identiques, avec quatre
candidats initialement au centre et des coûts initiaux égaux à 1.

Pour la graine maîtresse $r$ et le tour $t=0,\ldots,500$, la graine
32 bits est formée des quatre premiers octets, en petit-boutiste, du
SHA-256 de la chaîne UTF-8 \texttt{pouls:r:t}. Les tirages gaussiens
du moteur et les ex aequo des entrées sont ainsi reproductibles.
Les neuf séries sont stockées dans \texttt{simulation-results.json}.
La table et les figures sont générées depuis ces calculs, et non
recopiées d'une ancienne expérience de l'application.

\subsection{Participation et renouvellement}
% Generated by node simulation-results.cjs; do not edit.
\begin{center}\small
\begin{tabular}{lrrrr}
\hline
Sc\'enario & Graine & Abst.\ $t=500$ & Abst.\ $401$--$500$ & Entr\'ees\\
\hline
R\'ef\'erence & 314159 & 30.1\,\% & 26.4\,\% & 16\\
R\'ef\'erence & 271828 & 29.9\,\% & 27.0\,\% & 16\\
R\'ef\'erence & 161803 & 56.8\,\% & 53.1\,\% & 20\\
\hline
Sans r\'eputation & 314159 & 68.7\,\% & 69.6\,\% & 0\\
Sans r\'eputation & 271828 & 69.5\,\% & 69.8\,\% & 0\\
Sans r\'eputation & 161803 & 70.0\,\% & 69.6\,\% & 0\\
\hline
Quadratique & 314159 & 0.2\,\% & 0.6\,\% & 16\\
Quadratique & 271828 & 0.3\,\% & 0.9\,\% & 16\\
Quadratique & 161803 & 0.3\,\% & 0.7\,\% & 16\\
\hline
\end{tabular}\end{center}
\expandafter\def\csname abstMean0\endcsname{35.5}
\expandafter\def\csname abstMin0\endcsname{26.4}
\expandafter\def\csname abstMax0\endcsname{53.1}
\expandafter\def\csname abstMean1\endcsname{69.7}
\expandafter\def\csname abstMin1\endcsname{69.6}
\expandafter\def\csname abstMax1\endcsname{69.8}
\expandafter\def\csname abstMean2\endcsname{0.7}
\expandafter\def\csname abstMin2\endcsname{0.6}
\expandafter\def\csname abstMax2\endcsname{0.9}


Dans la référence, l'abstention moyenne sur les 100 derniers tours
vaut \csname abstMean0\endcsname\,\% en moyenne sur les trois graines ;
elle varie de \csname abstMin0\endcsname\,\% à
\csname abstMax0\endcsname\,\%. Deux trajectoires finissent autour
de 30\,\%, mais la troisième à 56,8\,\% : \textbf{30\,\% n'est pas
un niveau garanti par les paramètres par défaut}.
Les 16 à 20 nouvelles entrées montrent que les places restent actives
alors même que certains occupants disparaissent.

Sans réputation, l'abstention moyenne finale est proche de
\csname abstMean1\endcsname\,\%, sans renouvellement observé.
Dans ces trois essais, tous les candidats conservent une petite base
acceptable et leurs coûts restent à 1. Avec le coût quadratique,
la même statistique descend à \csname abstMean2\endcsname\,\%.
Cette comparaison conserve $\beta$ sans recalibrer l'échelle :
$d^2<d$ pour $0<d<1$, mais $d^2>d$ pour $d>1$.
Ce n'est donc pas une mesure d'un effet pur de courbure à couverture
initiale constante.

\subsection{Trajectoire illustrée et vidéo}
Les figures suivantes reprennent la référence, graine $314159$,
également utilisée pour la vidéo jointe. Au tour 500, les parts sur
tout l'électorat, classées, valent environ
$35{,}85$, $33{,}81$, $0{,}19$ et $0{,}09$\,\%.
La concentration observée ne reproduit donc pas une suite géométrique
imposée de type $40,20,10,5,\ldots$.
La vidéo accélérée parcourt les mêmes 501 états, sans recalculer une
trajectoire différente.
Pour la lecture imprimée, les figures ci-dessous conservent une couleur
de base par place ; l'application et la vidéo distinguent en plus les
partis successifs par leur teinte et leur numéro de génération.

\begin{figure}[htbp]\centering
\begin{tikzpicture}\begin{axis}[pouls,ymin=-1,ymax=1,ylabel={Position},
  legend columns=4,legend pos=north west]
\foreach \col in {A,B,C,D}{
  \addplot table[x=tour,y=p\col,col sep=comma]{simulation-trajectoires.csv};
  \expandafter\addlegendentry\expandafter{\col}
}
\end{axis}\end{tikzpicture}
\caption{Positions des quatre places de candidature, référence, graine 314159.
Les sauts peuvent correspondre à un renouvellement exempt de la borne $\gamma$.}
\end{figure}

\begin{figure}[htbp]\centering
\begin{tikzpicture}\begin{axis}[pouls,ymin=0,ylabel={Coût de base $k_i$}]
\foreach \col in {A,B,C,D}{
  \addplot table[x=tour,y=k\col,col sep=comma]{simulation-trajectoires.csv};
}
\end{axis}\end{tikzpicture}
\caption{Coûts de base appliqués. Une nouvelle entrée ramène le coût à 1 ;
les trajectoires ne suivent pas un unique candidat pendant 500 tours.}
\end{figure}

\begin{figure}[htbp]\centering
\begin{tikzpicture}\begin{axis}[pouls,ymin=0,ymax=100,
  ylabel={Part de l'électorat (\%)},legend columns=5,legend pos=north east]
\foreach \col in {A,B,C,D}{
  \addplot table[x=tour,y=v\col,col sep=comma]{simulation-trajectoires.csv};
  \expandafter\addlegendentry\expandafter{\col}
}
\addplot[black,dashed,thick] table[x=tour,y=abstention,col sep=comma]{simulation-trajectoires.csv};
\addlegendentry{Abstention}
\end{axis}\end{tikzpicture}
\caption{Voix espérées et abstention, rapportées à l'ensemble de l'électorat.}
\end{figure}
\clearpage

\subsection{Interprétation et limites}
La réputation produit un avantage cumulatif : un score élevé réduit le
coût, élargit l'admissibilité et peut renforcer le score suivant.
Le renouvellement réintroduit une offre à coût 1 auprès d'un groupe
précédemment délaissé, sans assurer sa survie.
Trois graines par scénario fournissent seulement une illustration
descriptive : ni intervalle de confiance, ni calibration empirique,
ni preuve de robustesse générale.

L'état dynamique comprend les positions, les coûts, les trois derniers
scores, les âges de mémoire des occupants et les tirages de coût.
Des positions momentanément immobiles ne prouvent donc pas un équilibre.
Le bruit renouvelé, les coûts évolutifs, les décisions simultanées et
les entrées rendent une convergence déterministe non garantie.
Les maxima sont ceux de grilles finies, pas une preuve d'optimalité
continue. Une étude complémentaire devrait varier les conditions
initiales, les graines, le nombre de candidats et la résolution, puis
tester les distributions non uniformes.

\paragraph{Reproduction.}
Exécuter ces commandes, puis lancer deux autres passes de compilation :
\begin{quote}\small
\begin{verbatim}
node simulation-results.cjs
pdflatex -jobname=pouls-modele modele-theorique.tex
bibtex pouls-modele
\end{verbatim}
\end{quote}
Le PDF produit est \texttt{pouls-modele.pdf}.
Les fichiers de données et la vidéo sont fournis à côté du document.

\bibliographystyle{plain}
\bibliography{references}
\end{document}
